\section{Metodología de simulación e implementación}
\subsection{Cálculo y simulación en MATLAB}
Se desarrolló un script de MATLAB que integra los tres niveles del análisis.
Primero calcula la referencia ideal con los estados
verdaderos, $u=-Kx+N_r r$, para aislar la dinámica de $A-BK$ sin observador;
luego reproduce ADC, DAC bipolar, saturaciones y cronograma del
firmware; finalmente importa las cuatro capturas UART/Rigol. Para cada caso
exporta pico y tiempo al pico, valor final, error RMS, esfuerzo máximo y RMS de
control, dispersión estacionaria del error observable, período de muestreo,
tiempo de CPU y margen temporal. También genera las trayectorias de estados, el
error observable, la convergencia ante una perturbación inicial y el análisis de
coincidencia de códigos DAC.

El error completo $\mathbf x-\widehat{\mathbf x}$ se evalúa en simulación. En
el banco solo $x_2=y$ posee una medición independiente; por ello la evidencia
experimental de error se expresa como $e_y=y-\hat x_2$, sin presentar
$x_1-\hat x_1$ como si hubiese sido medido.

\subsection{Arquitectura implementada en el PSoC}
En cada interrupción de muestreo se adquiere J7 mediante el ADC, se convierte a
unidades físicas y se actualiza el estimador seleccionado. Luego se calcula la
ley~\eqref{eq:control}, se limita la solicitud a la región validada y se escribe
el DAC. La UART registra referencia, salida, estados estimados, innovación,
mando solicitado, mando nominal, banderas, intervalo y tiempo de ejecución.

Los coeficientes identificados se integraron con límites conservadores de
$[-115,107]$ mV, detección abortiva de saturación y comprobación posterior a la
cuantización. Las cuatro combinaciones de arquitectura y ubicación compilaron
en PSoC Creator 4.4 y se programaron sucesivamente. Ninguna
ejecución generó banderas. Al finalizar se volvió a compilar y programar
MONITOR, con el DAC en cero y todas las autorizaciones de lazo cerrado
deshabilitadas; después se apagó la fuente bipolar.

\subsection{Protocolo de medición}
Antes del lazo cerrado se verificaron masas, sondas $10\times$, polaridad,
salida nula, respuesta abierta, período de muestreo y retorno seguro. Para cada
configuración se programó el firmware correspondiente, se rearmó el sistema y
se adquirieron simultáneamente 256 registros UART y 1200 puntos por canal del
Rigol. Los metadatos de cada adquisición permitieron verificar el modo, la
familia de polos, el modelo identificado, el disparo físico, la ausencia de
disparo forzado, $T_s$ y las banderas. Las ondas crudas se conservaron sin
modificaciones; el suavizado mostrado en una figura se utiliza únicamente para
visualización.
